\chapter{Introduction}
\label{ch:intro}

\section{Introduction}
Photographs and video of a person's face have long served as informal evidence of identity and events. Generative models have weakened
that assumption: generative adversarial networks (GANs) \citep{goodfellow2014gan,karras2019stylegan}, face-swapping and re-enactment
tools \citep{rossler2019ff} and, more recently, diffusion models \citep{ho2020ddpm,rombach2022ldm} can produce or alter faces whose
artefacts are invisible to the naked eye. Research on \emph{media forensics} \citep{verdoliva2020media,tolosana2020deepfakes} responds
with automatic detectors, but a detector that only prints ``fake'' or ``real'' is of limited use to a reviewer: the verdict has to be accompanied
by evidence that can be inspected, by the conditions under which it is reliable, and by an honest statement of what it cannot decide.

\emph{Verif-AI} is a software system built around that idea. It accepts an image or a video, finds the faces in it, scores every face with a
two-stream convolutional network, shows which facial regions influenced the score, summarises a video over time and generates a PDF
forensic report. This report describes the design, the implementation, and --- with equal weight --- the evaluation of the system,
including the experiments that showed where it fails.

\subsection{Problem Statement: Verifying Faces in an Age of Synthetic Media}
The practical problem has three parts.

\begin{enumerate}
\item \textbf{Detection.} Given a face crop, estimate whether it shows a photographed face or a generated or manipulated one. Detectors
trained on one family of manipulations often transfer poorly to another \citep{wang2020cnn,ojha2023universal,corvi2023diffusion}, so the
measured accuracy depends heavily on how the test data relate to the training data.
\item \textbf{Evidence.} A score alone does not let a reviewer judge a result. The system must localise what drove the decision, report the
detection confidence of the face finder, show where in a video suspicious frames occur and keep a record (file hash, model checkpoint,
settings) that makes the analysis reproducible.
\item \textbf{Honesty and privacy.} Neural scores are not probabilities of authenticity \citep{guo2017calibration}, detectors are exposed to
\emph{shortcuts} in their training data \citep{geirhos2020shortcut}, and uploaded media can be sensitive. The system must therefore say
``inconclusive'' when it cannot decide, avoid presenting scores as certainties, and delete uploads after use.
\end{enumerate}

\begin{figure}[ht]
\centering
\begin{adjustbox}{max width=\linewidth,center}
\begin{tikzpicture}[font=\small,node distance=0.45cm and 0.6cm,
  box/.style={draw,rounded corners=2pt,minimum height=0.9cm,align=center,fill=white},
  io/.style={draw,rounded corners=2pt,minimum height=0.9cm,align=center,fill=navy!8},
  arr/.style={-{Stealth[length=2mm]},thick}]
\node[io] (in) {Image or\\video};
\node[box,right=of in] (det) {Face\\detection\\(MTCNN)};
\node[box,right=of det] (model) {Two-stream\\classifier};
\node[io,right=of model] (out) {Verdict, score,\\Grad-CAM, report};
\draw[arr] (in)--(det); \draw[arr] (det)--(model); \draw[arr] (model)--(out);
\node[box,below=0.6cm of model,fill=accent!12] (ev) {Evidence: heat-maps, timeline, face tracks, PDF report};
\draw[arr,dashed] (model)--(ev);
\end{tikzpicture}
\end{adjustbox}
\caption[Conceptual overview of Verif-AI]{Conceptual overview of Verif-AI. A face is located, scored by a two-stream network and returned together with inspectable
evidence. Scores are raw model outputs (not probabilities); INCONCLUSIVE is an explicit outcome.}
\label{fig:overview}
\end{figure}

\subsection{Motivation and Project Objectives}
The project started from an earlier prototype that combined an RGB and an SRM noise stream and reported very high accuracy on a single
dataset. When the prototype's checkpoint was evaluated on independent, identity-disjoint data it performed at chance level, which made the
central motivation of this work concrete: \emph{a detector is only as trustworthy as its evaluation}. The objectives were therefore both
engineering and scientific:

\begin{itemize}
\item \textbf{Objective 1:} to re-implement the two-stream detector as a clean, tested software package that loads the original
checkpoints and a newly trained one, with identical preprocessing for training and inference.
\item \textbf{Objective 2:} to train and evaluate the model with identity-disjoint splits, duplicate removal, validation-only model
selection and a clear separation of internal and cross-generator evaluation.
\item \textbf{Objective 3:} to provide explainable output --- Grad-CAM heat-maps, face-level scores, video timelines, face tracks and a PDF report --- through an
API and a web dashboard, with privacy-preserving handling of uploads.
\item \textbf{Objective 4:} to investigate honestly where the model fails (for example on genuine video) and to document the cause and the
limits of every claim.
\item \textbf{Objective 5:} to deploy a public demonstration (static frontend, hosted API) and to state its restrictions truthfully.
\end{itemize}

\subsection{Scope and Target Applications}
Verif-AI analyses \emph{faces}. It is intended for educational demonstration, exploratory review of suspicious face media by a human and
research on detector generalisation. It is \textbf{not} a general detector of AI-generated images or videos, does not analyse non-face content,
does not establish authorship or provenance, and a ``real'' result is not proof of authenticity. Scores are uncalibrated; these boundaries decide how
every number in Chapter~\ref{ch:results} may be read.

\subsection{Contributions}
\begin{enumerate}
\item A tested two-stream detection package with a documented preprocessing pipeline that reproduces training-time cropping, JPEG
round-trip and normalisation at inference time (Chapters~\ref{ch:design} and~\ref{ch:impl}).
\item An image and video analysis service with Grad-CAM evidence, face tracking, frame timelines, PDF reports and deletion
(Chapter~\ref{ch:impl}), and a dashboard with original, heat-map and overlay views.
\item An evaluation record that includes an identified failure (chance-level behaviour on genuine video), a diagnosis of its cause as an
image-quality shortcut, a pre-registered acceptance gate and modern-generator experiments (Chapter~\ref{ch:results}).
\item A reproducible training and data-review workflow (Kaggle GPU notebooks that export only checkpoints and metrics) and a dataset
licence review that did not accept licence badges as proof (Chapter~\ref{ch:results} and Appendix~\ref{app:datasets}).
\end{enumerate}

\subsection{Organisation of the Report}
Chapter~\ref{ch:lit} surveys the field, the research gap and the contributions of this work; Chapter~\ref{ch:req} the requirements and technology stack;
Chapter~\ref{ch:design} the architecture, model and pipelines; Chapter~\ref{ch:impl} their implementation, testing and deployment;
Chapter~\ref{ch:results} the experiments, their interpretation and limitations; Chapter~\ref{ch:concl} the conclusion and future work.
Appendices~\ref{app:code}--\ref{app:datasets} give key source code, the API and configuration, and the dataset review.
